\subsection{Related Work on Global Robustness}

%
\eca{Some works rely on specialized architectures or dedicated robust training procedures that ofer better global bounds, but suffer from lower clean accuracy \citep{Leino,Zhang22a,Sun22,Singla22, Chen21}. However, these methods are also frequently constrained by specific thread models: \citet{Singla22} is restricted to the $L_2$ norm, \citet{Zhang22a} to the $L_{\infty}$ norm, while \citet{Chen21} considers non-$L_p$ security properties. \citet{REGLO} provides a post-hoc DNN repair method and is thus restricted to local/in-distribution (ID) regions rather than the full-input domain.
In contrast, our approach postprocesses standard pretrained DNNs by integrating (linear) autoencoders to filter out non-ID features. This allows to tightly control the accuracy loss by tuning the latent dimension size.}

\iffalse
Some works focus on training DNNs to be globally robust. Specifically, \cite{Leino,Zhang22a,Sun22,Singla22} propose special architectures with better global bounds, but lower accuracy.
Further, \citet{Leino,Chen21,Singla22} do not consider $L_\infty$ or $L_1$ norm, while \citet{REGLO} do not consider (non-ID) inputs.
%A key aspect of global robustness involves ensuring that analyzed data is both semantically meaningful and in-distribution (ID). 
Our approach is to postprocess DNNs to obtain better global bounds using (linear) auto encoders
%The (L)AE-DNN architectures (Fig.~\ref{fig:AE-DNN} in Section \ref{sec.pca}) project the input on a latent space 
to filter out non-ID features, with limited loss of accuracy, and can be applied as postprocessing to these DNNs.
\fi

Some works \citep{Bastani16,Ruan19,Gopinath18} estimate global robustness bounds, but without hard verification certificates, e.g., using probabilistic guarantees \citep{Levy23,Mangal19}.
%More related, s
Others \citep{Marabou,lipshitz,ITNE,GROCET} consider certifying restricted Lipschitz bound, to understand how the output can change given an input perturbation. However, they do not consider whether the classification changes. 
%
Closest to our work is VHAGaR \citep{vhagar}, which computes 
bounds $\bar{\alpha}^\varepsilon_{C,D}$ similar to our bounds 
$\bar{\beta}^\varepsilon_{C,D}$.
%, but limited to $L_\infty$-perturbations and variants. 
Crucially, VHAGaR bounds are not small enough for real-time robustness, whereas our bounds are. 
%One reason is {\em Diff} MILP is more accurate in practice.

Technically, while \cite{diff} does not tackle global robustness, it is one of the first papers to consider {\em differential} verification, introducing the {\em diff} variables that are used here and partly in \cite{ITNE}. 
However, \citet{diff} did not consider MILP encoding.
The MILP-based works \citep{vhagar,lipshitz,ITNE} 
use the classical MILP encoding from \cite{MILP}.
%, and not our novel {\em Diff} model, nor even the "1b" model. 
ITNE \citep{ITNE} adds explicitly one {\em linear} constraint (Eq. (3) in \cite{ITNE}) on the {\em diff variables} $y,\hat{y}$,
%$\frac{y_i-\gamma_i}{2} \leq \hat{y}_i \leq \frac{y_i+\gamma_i}{2}$ 
implied by each of our MILP models {\em Diff}, {\em 2b+1} and {\em 1b} models. VHAGaR \citep{vhagar} instead adds per neuron constraints (depending on the perturbation) between the (binary) variables to simplify the MILP problem, which is well-adapted to occlusion and patches but less to $L_1$-perturbations.
None of these works consider $L_1$-perturbations.



%While non-linear AEs can achieve better compression and noise reduction, we consider linear AE in this work to avoid the additional computational complexity associated with non-linear activation functions within the MILP framework~\cite{fischetti2018deep}.

